\documentclass[11pt,letterpaper]{article}
\usepackage[T1]{fontenc}
\usepackage[utf8]{inputenc}
\usepackage{lmodern,microtype}
\usepackage[margin=0.79in,headheight=15pt]{geometry}
\usepackage{parskip,enumitem,longtable,tabularx,booktabs,array,xcolor,fancyhdr,hyperref}
\definecolor{navy}{HTML}{17314C}
\definecolor{teal}{HTML}{087F78}
\hypersetup{colorlinks=true,linkcolor=navy,urlcolor=teal}
\setlength{\emergencystretch}{3em}
\setlist[itemize]{leftmargin=*,nosep,topsep=5pt}
\setlist[enumerate]{leftmargin=*,itemsep=3pt}
\renewcommand{\arraystretch}{1.22}
\pagestyle{fancy}\fancyhf{}
\fancyhead[L]{\small Enterprise Data Modeling}\fancyhead[R]{\small Documentation Suite}
\fancyfoot[C]{\thepage}
\newcommand{\srcbook}{\textit{The Data Model Resource Book}, revised edition, Volume 2}
\newcommand{\ResponseLine}{\par\vspace{1.05em}\noindent\rule{\linewidth}{0.4pt}\par}
\newcommand{\ResponseArea}{\ResponseLine\ResponseLine\ResponseLine}
\begin{document}
\begin{center}
{\LARGE\bfseries\color{navy} Enterprise Data Modeling Curriculum}\par
\vspace{0.25em}{\large Twelve-week prerequisite-based study program}\par
\vspace{0.3em}{\small Original engineering study material \textbullet{} 2026 edition}
\end{center}
\vspace{0.8em}
\section{Learning objectives}
Build competence from conceptual modeling through verified physical implementation and governed enterprise adoption. The same evidence chain is used throughout: requirement → semantic rule → diagram → DDL constraint or application invariant → positive/negative test → architectural approval.


\section{Prerequisites and 12-week roadmap}
\begin{tabularx}{\linewidth}{@{}p{0.07\linewidth} p{0.22\linewidth} p{0.38\linewidth} p{0.29\linewidth}@{}}
\toprule
\textbf{Week} & \textbf{Topic} & \textbf{Required input} & \textbf{Exit artifact} \\ 
\midrule
1 & Prerequisites & Keys, sets, relational algebra, business glossary & 15-term glossary; sample relations \\ 
2 & Modeling fundamentals & Ch. 1; entity identity, cardinality, subtypes & Conceptual ERD, cardinality audit \\ 
3 & Logical normalization & Functional dependencies, 1NF to BCNF & Dependency proof and decomposition \\ 
4 & Physical PostgreSQL & PK, FK, unique, check, transaction semantics & DDL with smoke tests \\ 
5 & Manufacturing & Ch. 2; BOM, revisioning, work effort & Manufacturing ERD and BOM rule tests \\ 
6 & Telecom / Health Care & Ch. 3--4; network and encounters & UML diagrams and privacy checklist \\ 
7 & Insurance / Finance & Ch. 5--6; coverage and ledger & Temporal contract and balancing tests \\ 
8 & Services / Travel & Ch. 7--8; engagement and itinerary & Fulfillment/revision workflow models \\ 
9 & E-Commerce & Ch. 9; offers, orders, payment & Marketplace case with checkout DDL \\ 
10 & Cross-industry adaptation & Ch. 10; fit-gap and semantic ownership & Canonical-vs-context decision matrix \\ 
11 & Operations and governance & Classification, lineage, retention, migration & Governance review and expand-contract plan \\ 
12 & Capstone & Design defense, adverse cases, performance & Source-linked architecture assessment \\ 
\bottomrule
\end{tabularx}

\section{Prerequisite gates}
\begin{itemize}
\item Gate A: Can identify candidate keys and derive join multiplicity from requirements.
\item Gate B: Can state functional dependencies and exhibit a lossless decomposition.
\item Gate C: Can run PostgreSQL DDL and interpret expected constraint violations.
\item Gate D: Can explain valid time, system recorded time and how migration affects consumers.
\item Gate E: Can assess conceptual/logical/physical consistency with traceable evidence.
\end{itemize}

\section{Learning sequencing}
Do not start industry data-model implementation until the Week 1--4 foundations are complete. The financial, health care and insurance tracks add regulated-data considerations and should be reviewed by domain stakeholders before implementation. Weeks 5--9 may be distributed across a team only after all participants pass the foundations gate.


\section{Capstone exit criteria}
\begin{tabularx}{\linewidth}{@{}p{0.28\linewidth} p{0.68\linewidth}@{}}
\toprule
\textbf{Criterion} & \textbf{Minimum evidence} \\ 
\midrule
Model correctness & No unexplained many-to-many relation, missing key or contradictory optionality \\ 
Normalization & Documented dependencies with justified exception for denormalization \\ 
Executable behavior & At least eight test cases including five negative/boundary paths \\ 
Governance & Owner, classification, retention and lineage for critical entities \\ 
Migration readiness & Versioned forward deployment, backfill, verification and rollback decision \\ 
Architecture assessment & Scorecard >= 80/100 with no critical security or integrity blocker \\ 
\bottomrule
\end{tabularx}

\vfill
\hrule
{\footnotesize\textbf{Provenance.} Supplementary engineering practice: original content; the cited book provides domain-modeling context.}
\end{document}
